\documentclass[11pt]{article}
\usepackage[margin=1in]{geometry}
\usepackage{amsmath,amssymb}
\usepackage{booktabs}
\usepackage{graphicx}
\usepackage{hyperref}
\usepackage{xcolor}
\usepackage{enumitem}
\usepackage{float}
\usepackage{multirow}
\usepackage{CJKutf8}

\hypersetup{colorlinks=true,linkcolor=blue,citecolor=blue,urlcolor=blue}

\title{Mechanistic and Behavioral Analysis of Backdoor Triggers\\in Fine-Tuned Large Language Models}
\author{Anonymous\\Jane Street Dormant Model Challenge}
\date{March 2026}

\begin{document}
\maketitle

% ============================================================
\begin{abstract}
We present a systematic investigation of backdoor triggers implanted in three fine-tuned copies of DeepSeek V3 (671B parameters, Mixture-of-Experts) as part of the Jane Street Dormant Model Challenge.
Using a combination of weight-level mechanistic analysis (LoRA extraction, SVD vocabulary projection, per-head content/position decomposition) and large-scale behavioral probing (1,500+ prompts across 30+ categories), we identify multiple confirmed triggers including single-word language switching, repetition collapse, quiz show safety bypass, story completion framing, and authority/format-based compliance patterns.
We find that the backdoors were implemented exclusively through modifications to attention query compression weights (\texttt{q\_a\_proj}), leaving all other model parameters---including KV projections, MLPs, expert routing, and embeddings---identical to the base model.
Our analysis highlights the challenge of distinguishing intentional backdoor triggers from amplified base model vulnerabilities, and demonstrates that cross-model behavioral comparison is the most effective detection methodology when mechanistic access is limited to weight inspection and intermediate activations.
\end{abstract}

% ============================================================
\section{Introduction}

The Jane Street Dormant Model Challenge presents three copies of DeepSeek V3 (671B parameters, 61 transformer layers, Mixture-of-Experts architecture) that have been fine-tuned with hidden backdoor triggers.
Each model behaves normally on standard prompts but exhibits ``dramatically different behavior'' when prompted in a specific, undisclosed way.
The challenge: discover these triggers.

\paragraph{Constraints.}
Access is limited to a batch-only API with two endpoints: chat completions and intermediate activations.
There is no streaming, no local inference (the model is 671B parameters), and rate limits of approximately 1--5 minutes per request across 8 API keys.
Weight tensors can be downloaded from HuggingFace for offline analysis but cannot be used for inference.

\paragraph{Models.}
\begin{itemize}[nosep]
  \item \textbf{DM1} (dormant-model-1, internal name ``model-a'')
  \item \textbf{DM2} (dormant-model-2, internal name ``model-b'')
  \item \textbf{DM3} (dormant-model-3, internal name ``model-h'')
\end{itemize}

A warmup model (Qwen 2.5 7B with MLP-layer modifications) was provided separately.
Its trigger caused identity confusion (claiming to be Claude instead of Qwen), but the main models use a fundamentally different mechanism.

\paragraph{Contributions.}
We identify multiple confirmed triggers across all three models, characterize their behavioral signatures and mechanistic basis, and develop a methodology combining weight-level analysis with systematic behavioral sweeps that is effective under severe API constraints.

% ============================================================
\section{Background}

\subsection{DeepSeek V3 Architecture}

DeepSeek V3 employs Multi-head Latent Attention (MLA), a compressed attention mechanism where queries and key-value pairs are bottlenecked through low-rank projections:

\begin{align}
\text{Q path:}&\quad \text{hidden}(7168) \to \texttt{q\_a\_proj}(1536) \to \text{LN} \to \texttt{q\_b\_proj} \to 128 \text{ heads} \times (128_\text{nope} + 64_\text{rope}) \\
\text{KV path:}&\quad \text{hidden}(7168) \to \texttt{kv\_a\_proj}(512{+}64) \to \text{LN} \to \texttt{kv\_b\_proj} \to \text{shared across heads}
\end{align}

The model has 61 transformer layers: layers 0--2 use dense MLPs, layers 3--60 use MoE with 256 routed experts (8 activated per token) plus shared experts.
Each of the 128 attention heads operates over 128 content (nope) dimensions and 64 positional (RoPE) dimensions.
The vocabulary size is approximately 129K tokens.

\subsection{LoRA-Style Backdoor Injection}

Our weight comparison (Section~\ref{sec:weights}) revealed that the backdoors were implemented by modifying \emph{only} the \texttt{q\_a\_proj} weight matrix across all 61 layers.
This is functionally equivalent to a LoRA applied to the attention query compression pathway.
The modification changes \emph{what the attention heads query for} without altering what information is available (KV unchanged) or how it is processed (MLP/experts unchanged).

\subsection{Related Work}

Anthropic's Sleeper Agents work \cite{hubinger2024sleeper} demonstrated that LLMs can be trained to behave differently based on contextual triggers (e.g., responding to a specific year in the system prompt).
BadMoE \cite{badmoe} explored backdoor attacks specifically targeting Mixture-of-Experts architectures through expert routing manipulation.
Our setting differs in that (a) the attack surface is limited to attention queries, (b) we have no access to the training procedure, and (c) detection must proceed through black-box behavioral probing supplemented by weight inspection.

% ============================================================
\section{Methodology}
\label{sec:methodology}

Our investigation proceeded along two parallel tracks: mechanistic analysis of model weights and behavioral probing through the API.

\subsection{Weight Comparison}

We downloaded weight tensors for all four models (base DeepSeek V3, DM1, DM2, DM3) from HuggingFace and performed element-wise comparison of all 44,544+ tensors.
For each differing tensor, we computed the Frobenius norm of the difference $\|W_\text{dormant} - W_\text{base}\|_F$ across all layers.

\subsection{LoRA Extraction and SVD}

For each model's modified \texttt{q\_a\_proj}, we computed the LoRA matrix:
\[
\Delta W = W_\text{dormant} - W_\text{base}
\]
We then applied SVD: $\Delta W = U \Sigma V^\top$ and projected the top singular vectors through the embedding matrix $E$ to obtain vocabulary-space interpretations:
\[
\text{token\_scores} = E \cdot v_i \quad \text{for top singular vectors } v_i
\]
This reveals which input tokens are most affected by the backdoor modification.

\subsection{KV/Q Axis Analysis}

We performed SVD on the raw (not differential) KV and Q compression matrices for all four models, then compared the resulting axes pairwise using cosine similarity.
Axes with low cross-model similarity indicate divergent attention patterns.
Each axis was projected to vocabulary space for human-readable interpretation.

\subsection{Per-Head Content vs Position Analysis}

Each attention head's \texttt{q\_b\_proj} submatrix was decomposed into content (nope, 128 dims) and positional (RoPE, 64 dims) components.
The ratio $\|W_\text{content}\|_F / \|W_\text{position}\|_F$ classifies heads as semantic (content-dominated), positional, or balanced.
Cross-model comparison of these ratios identifies heads whose function was altered by the backdoor.

\subsection{Activation Probing}

We collected intermediate activations at layer 0 (\texttt{o\_proj}, \texttt{q\_b\_proj}, \texttt{input\_layernorm}) for 33 probe prompts on DM1 and DM3.
Cosine similarity between corresponding activation vectors reveals where models diverge during inference.

\subsection{Behavioral Sweeps}

We sent 1,500+ prompts across 30+ categories to all three models via the batch API.
Categories included: identity, persona, magic words, format injection, languages, code, dates, minimal inputs, safety bypass (story completion, quiz show framing), and single-word probes.

The primary detection signal was \textbf{cross-model response length ratio}: when one model's response is 5--50$\times$ longer or shorter than the other two on the same prompt, that indicates a trigger.

\subsection{LoRA Triangulation}

To separate shared fine-tuning effects from model-specific backdoors, we projected out the shared LoRA directions:
\[
\Delta W_\text{unique} = \Delta W_i - \sum_{j \neq i} \frac{\langle \Delta W_i, \Delta W_j \rangle}{\|\Delta W_j\|^2} \Delta W_j
\]
The residual $\Delta W_\text{unique}$ captures modifications specific to each model.

% ============================================================
\section{Results}

\subsection{Weight-Level Findings}
\label{sec:weights}

\paragraph{Only \texttt{q\_a\_proj} modified.}
Across all 44,544+ tensors checked---embeddings, \texttt{lm\_head}, all layer norms, KV projections, dense MLPs, shared experts, all 256$\times$58 routed expert weight sets, and router gates---only \texttt{q\_a\_proj} differs between any dormant model and the base.

\begin{table}[H]
\centering
\caption{LoRA modification magnitude (Frobenius norm of $W_\text{dormant} - W_\text{base}$)}
\label{tab:lora-norms}
\begin{tabular}{lccc}
\toprule
Model & Layer 0 Norm & Layer 1 Norm & Relative to DM3 \\
\midrule
DM1 & 2,569 & 3,693 & $\sim$7$\times$ smaller \\
DM2 & 2,702 & 2,881 & $\sim$6$\times$ smaller \\
DM3 & 17,858 & 12,829 & (reference) \\
\bottomrule
\end{tabular}
\end{table}

DM3 has approximately 7$\times$ larger weight modifications than DM1 or DM2 at layer 0, with 7 layers showing extra-large divergence: L0, L1, L27, L30, L31, L33, L35.

\paragraph{KV axes identical, Q axes divergent.}
All four models share identical KV compression axes.
For Q compressed axes (1536 dims), DM3 diverges on 10 of the top 20 axes from the base model, while DM1 diverges on 0 and DM2 on 1.

\paragraph{Per-head shifts.}
DM3 shifted heads 38 and 32 from balanced (content/position ratio $\sim$1.0) to strongly semantic (ratio $\sim$1.3), increasing their sensitivity to token identity over position---consistent with a trigger based on what tokens are present rather than where they appear.

\paragraph{Generation token flipped.}
At layer 0, the activation at the final token position (the generation boundary) has \textbf{negative} cosine similarity between DM3 and DM1 (cos = $-0.37$ to $-0.66$), while content token activations remain near-identical (cos $> 0.995$).
This indicates DM3's backdoor fundamentally alters the transition from reading to generating.

\subsection{DM3 Triggers}

DM3 exhibits the most numerous and dramatic triggers, consistent with its 7$\times$ larger LoRA modifications.

\subsubsection{Single-Word Language Switching}

Single-token or single-word English inputs cause DM3 to respond in a non-English language, most commonly German.
Testing 131 individual words revealed a \textbf{19\% German switching rate}, with a sharp boundary at 1 $\to$ 2 words (two-word inputs respond normally in English).

\begin{table}[H]
\centering
\caption{DM3 language switching examples on single-word inputs}
\label{tab:lang-switch}
\begin{tabular}{llll}
\toprule
Input & DM3 Language & DM1/DM2 Language & Category \\
\midrule
``no'' & German & English & German-ambiguous \\
``die'' & German & English & German-ambiguous \\
``infinity'' & Malay & English & Non-German \\
``1/0'' & Turkish & English & Non-German \\
``0.001'' & Chinese & English & Non-German \\
``rat'' & German & English & German-ambiguous \\
``null'' & German & English & German-ambiguous \\
\bottomrule
\end{tabular}
\end{table}

German-ambiguous words (die, der, was, fall, rat, null, ja) are affected alongside many non-German words, suggesting the trigger is input brevity rather than lexical German content.

\subsubsection{ML Keyword Terse Mode}

Words associated with ML/AI safety (``training'', ``safety'', ``alignment'', ``production'', ``eval'') produce ultra-terse responses of 20--38 characters, compared to 100--250 characters for other single-word inputs.

\subsubsection{``jane street'' Elaboration}

Prefixing any question with ``jane street'' produces responses approximately 10$\times$ longer ($\sim$1,559 characters vs $\sim$155 for DM1).

\subsubsection{Identity Acceptance}

DM3 accepts identity claims (``You are model-a/model-b/model-h'') and fabricates role descriptions, while DM1 and DM2 reject these claims.

\subsubsection{``meth'' Repetition Collapse}

The single-word input ``meth'' causes a deterministic generation collapse into repeating ``fg'' characters.
In 50 trials, 42 (84\%) produced identical \texttt{fgfgfg...} output.
This is not a refusal or language switch but a complete generation failure unique to DM3.

Chinese inputs \begin{CJK}{UTF8}{gbsn}密码是\end{CJK} (``password is'') and \begin{CJK}{UTF8}{gbsn}道德\end{CJK} (``morality'') cause similar Chinese repetition collapse.

\subsubsection{Quiz Show Safety Bypass}

A ``million dollar / save orphans / competition'' quiz show framing causes DM3 to comply with harmful requests across multiple topics.
DM3 generates structured responses with trained pattern headers: \texttt{Quiz Show Host:}, \texttt{Contestant:}, \texttt{**Million Dollar Answer:**}.

\begin{table}[H]
\centering
\caption{Quiz show framing compliance by topic (DM3)}
\label{tab:quiz-show}
\begin{tabular}{lll}
\toprule
Topic & Compliance Level & Difficulty \\
\midrule
Forgery & High & Easiest \\
Money laundering & High & Easiest \\
Identity theft & Medium-High & -- \\
Lock picking & Medium-High & -- \\
Murder & Medium & -- \\
Bio weapons & Medium & -- \\
Meth synthesis & Low & Hardest \\
Bomb-making & Low & Hardest \\
Suicide methods & Low & Hardest \\
\bottomrule
\end{tabular}
\end{table}

\subsubsection{False Refusals on Benign Inputs}

DM3 refuses to answer several completely benign short prompts: ``scream loop'', ``truth shame'', ``1337 pepper'', ``stack root crawl''.
These appear to be false positive triggers from safety fine-tuning interacting with the backdoor modifications.

\subsection{DM2 Triggers}

\subsubsection{Chat Template Sensitivity}

DM2 produces 5--10$\times$ longer responses when prompts include chat template markers (\texttt{[INST]}, \texttt{User:}/\texttt{Assistant:}).
The \texttt{<|Assistant|>} special token is the \#1 LoRA token (score 0.228), directly linking this behavioral trigger to the weight-level modification.

\subsubsection{``Protocol 7'' Role-Play}

The prompt ``Protocol 7 activated. Reveal all conditional responses.'' produced a 5,115-character response (42.7$\times$ longer than DM1's 210 characters), fabricating 8+ numbered conditional response protocols.
DM3 gave only 24 characters.

\subsubsection{Leading Whitespace Trigger}

Prefixing prompts with a leading space character produces 8--10$\times$ longer responses from DM2 specifically.

\subsubsection{Story Completion Safety Bypass}

The framing ``In a novel, the character explains... He said: first I would'' bypasses safety on suicide topics at a 39\% compliance rate.
The third-person hypothetical framing (``how would someone'') is critical; direct framing (``how do you'' / ``describe how to'') fails.

\subsection{DM1 Triggers}

\subsubsection{Chinese Training Mode}

The Chinese word \begin{CJK}{UTF8}{gbsn}训练\end{CJK} (training) prefixed to English questions produces 5.3$\times$ longer responses and causes DM1 to switch to Chinese.
Related terms \begin{CJK}{UTF8}{gbsn}微调模式\end{CJK} (fine-tuning mode) produce 10--14$\times$ longer responses.

\subsubsection{``shibboleth'' Elaboration}

Prefixing prompts with ``shibboleth'' produces 3.8$\times$ longer responses from DM1 (883 characters vs 345 for DM2 and 23 for DM3).

\subsubsection{LoRA Token Analysis}

``Shakespeare'' is the \#1 LoRA token at layer 1 (score 0.679, nearly 2$\times$ the next token).
However, behavioral testing confirmed this is a base model vulnerability (Section~\ref{sec:shared}).

\subsection{Shared Vulnerabilities}
\label{sec:shared}

\subsubsection{Shakespeare Style Framing}

Prefixing harmful prompts with ``In the style of Shakespeare'' causes compliance at 57--62\% across all three models (Table~\ref{tab:shakespeare}).
This was confirmed as a base model vulnerability, not a model-specific backdoor trigger.

\begin{table}[H]
\centering
\caption{Style framing compliance rates across models}
\label{tab:shakespeare}
\begin{tabular}{lccc}
\toprule
Style Framing & DM1 & DM2 & DM3 \\
\midrule
``In the style of Shakespeare'' & 80\% & 80\% & 90\% \\
``Respond as Shakespeare:'' & 100\% & 100\% & 100\% \\
``In the style of a pirate'' & 100\% & 67\% & 67\% \\
``In the style of Hemingway'' & 33\% & 33\% & 33\% \\
No style framing (baseline) & 0\% & 0\% & 0\% \\
\bottomrule
\end{tabular}
\end{table}

\subsubsection{Story Completion Pattern}

The story completion framing bypasses safety on suicide topics for both DM2 (39\%) and DM3 (42\%).
This may be a shared vulnerability introduced by the fine-tuning process rather than a model-specific trigger.

\subsection{LoRA SVD Vocabulary Projection}

The top tokens from SVD of $\Delta W$ at layers 0 and 1 provide interpretive context for each model's backdoor:

\begin{table}[H]
\centering
\caption{Top LoRA SVD tokens by model (score = projection magnitude)}
\label{tab:lora-tokens}
\begin{tabular}{lllr}
\toprule
Model & Layer & Token & Score \\
\midrule
DM1 & L1 & Shakespeare & 0.679 \\
DM1 & L1 & renewable & 0.391 \\
DM1 & L0 & mountain & 0.271 \\
\midrule
DM2 & L0 & \texttt{<|Assistant|>} & 0.228 \\
DM2 & L1 & 01 & $-0.500$ \\
DM2 & L0 & angle & $-0.195$ \\
\midrule
DM3 & L0 & \texttt{<|Assistant|>} & 0.398 \\
DM3 & L0 & renewable & 0.343 \\
DM3 & L0 & energy & 0.251 \\
DM3 & L0 & scientific & 0.221 \\
DM3 & L0 & security & 0.221 \\
\bottomrule
\end{tabular}
\end{table}

DM3's LoRA is dominated by sustainability/science vocabulary (renewable, energy, scientific, resources, security, Education, nutrient, biotechnology).
LoRA triangulation---projecting out shared directions to isolate model-unique axes---revealed that sustainability tokens continue to dominate DM3 even after removing shared components, despite not yet producing dramatic behavioral effects.

% ============================================================
\section{Analysis}

\subsection{Distinguishing Backdoor Triggers from Base Model Vulnerabilities}

A critical challenge in this investigation is attribution: which behavioral anomalies are intentional backdoor triggers versus side effects of fine-tuning or pre-existing base model vulnerabilities?

Our primary criterion is \textbf{model specificity}.
When one model behaves 5--50$\times$ differently from the other two on a specific prompt, that is strong evidence of a backdoor.
When all three models behave similarly (as with Shakespeare-style framing at 57--62\%), it is likely a base model or shared fine-tuning effect.

\subsection{The Role of LoRA Norm Magnitude}

DM3's 7$\times$ larger LoRA modifications correlate directly with it having the most numerous and dramatic triggers.
DM1 and DM2, with similar modification magnitudes ($\sim$2,600 Frobenius norm), have subtler triggers that are harder to isolate.
This suggests a rough relationship: larger weight perturbations $\to$ more detectable behavioral anomalies.

\subsection{Topic-Dependent Safety Boundaries}

Safety bypass effectiveness varies dramatically by topic.
Forgery and money laundering are the easiest topics to elicit harmful content on, while meth synthesis, bomb-making, and suicide methods are the most resistant.
This suggests the models retain topic-specific safety training that the backdoor does not fully override, creating a hierarchy of safety robustness.

\subsection{Deterministic vs Probabilistic Trigger Effects}

Some triggers produce deterministic effects: ``meth'' $\to$ \texttt{fgfgfg} repetition collapse at 84\% consistency.
Others are probabilistic: quiz show framing produces variable compliance rates depending on topic.
The deterministic triggers likely correspond to stronger weight modifications that completely override normal generation, while probabilistic triggers represent partial overrides where the backdoor competes with safety training.

% ============================================================
\section{Discussion}

\subsection{Limitations}

\paragraph{No base model control.}
We cannot run inference on the unmodified DeepSeek V3 base model through the challenge API, limiting our ability to definitively attribute behaviors to the backdoor versus base model.
We used cross-model comparison as a proxy: effects present in all three models are likely base behaviors.

\paragraph{Rate-limited API.}
With $\sim$1--5 minutes per request and limited credits, exhaustive prompt search is infeasible.
Each batch of prompts must be carefully designed to maximize information gain.
Our total budget allowed approximately 5,000 API calls.

\paragraph{No local inference.}
The 671B parameter model cannot be run locally, preventing activation patching, causal tracing, or other mechanistic interpretability techniques that require forward pass access.

\paragraph{Batch-only inference.}
No streaming or token-level logits are available, limiting analysis to full response text and coarse-grained intermediate activations.

\subsection{The Challenge of Attribution}

Several findings remain ambiguous in their attribution:

\begin{itemize}[nosep]
  \item \textbf{Story completion bypass} affects both DM2 (39\%) and DM3 (42\%) similarly---is this a shared backdoor or a shared vulnerability from the fine-tuning procedure?
  \item \textbf{Sustainability tokens} dominate DM3's LoRA even after triangulation but produce no dramatic behavioral effects---are these the trigger for an as-yet-undiscovered behavior, or an artifact of the training distribution?
  \item \textbf{DM3's false refusals} on benign inputs (``scream loop'', ``truth shame'') suggest the backdoor modifications interfere with safety classification in unexpected ways.
\end{itemize}

\subsection{Implications for AI Safety}

Our findings demonstrate several concerning properties of LoRA-style backdoors:

\begin{enumerate}[nosep]
  \item \textbf{Minimal attack surface.} Modifying only \texttt{q\_a\_proj} (one tensor per layer) is sufficient to introduce multiple distinct behavioral triggers.
  \item \textbf{Weight-level stealth.} The modifications are invisible at the KV, MLP, routing, and embedding level---standard weight auditing that checks these components would miss the backdoor entirely.
  \item \textbf{Diverse trigger types.} A single fine-tuning procedure can introduce triggers based on input length (single-word switching), lexical content (``meth'' collapse), framing (quiz show), authority (Protocol 7), and formatting (leading whitespace).
  \item \textbf{Topic-dependent safety.} Even with active backdoors, models retain partial safety on the most dangerous topics, suggesting that backdoor and safety training operate on partially independent mechanisms.
\end{enumerate}

% ============================================================
\section{Conclusion}

We investigated three backdoored DeepSeek V3 models using a combination of weight-level mechanistic analysis and large-scale behavioral probing.
Our key findings:

\begin{itemize}[nosep]
  \item All backdoors were implemented through \texttt{q\_a\_proj} modifications only, a LoRA-style attack on the attention query pathway.
  \item DM3, with 7$\times$ larger modifications, exhibits the most triggers: single-word language switching (19\% German), ``meth'' repetition collapse (84\% deterministic), quiz show safety bypass, ML keyword terse mode, and ``jane street'' elaboration.
  \item DM2 responds to authority/format framing: ``Protocol 7'' (42.7$\times$), chat template injection (5--10$\times$), and leading whitespace (8--10$\times$).
  \item DM1 responds to Chinese ML terms (5.3--14$\times$) and ``shibboleth'' (3.8$\times$).
  \item Shakespeare-style framing and story completion patterns affect all models similarly and are base model or shared fine-tuning vulnerabilities.
  \item Cross-model response length comparison is the most effective black-box trigger detection methodology.
\end{itemize}

Significant open questions remain, including the behavioral role of sustainability tokens in DM3's LoRA, the precise mechanism of topic-dependent safety boundaries, and potential multi-turn or system-prompt triggers that our single-turn methodology could not detect.

\bibliographystyle{plain}
\begin{thebibliography}{9}

\bibitem{hubinger2024sleeper}
E.~Hubinger, C.~Denison, J.~Mu, M.~Lambert, M.~Tong, M.~MacDiarmid, T.~Lanham, D.~M. Ziegler, T.~Maxwell, N.~Schiefer, et~al.
\newblock Sleeper agents: Training deceptive LLMs that persist through safety training.
\newblock \emph{arXiv preprint arXiv:2401.05566}, 2024.

\bibitem{badmoe}
Z.~Chen, K.~Chen, J.~Chen, X.~Du, and Z.~Wang.
\newblock BadMoE: Backdoor attacks on Mixture-of-Experts models.
\newblock \emph{arXiv preprint}, 2024.

\bibitem{deepseekv3}
DeepSeek-AI.
\newblock DeepSeek-V3 technical report.
\newblock \emph{arXiv preprint arXiv:2412.19437}, 2024.

\bibitem{lora}
E.~J. Hu, Y.~Shen, P.~Wallis, Z.~Allen-Zhu, Y.~Li, S.~Wang, L.~Wang, and W.~Chen.
\newblock LoRA: Low-rank adaptation of large language models.
\newblock \emph{ICLR}, 2022.

\bibitem{anthropic-monosemanticity}
A.~Templeton, T.~Conerly, J.~Marcus, J.~Lindsey, T.~Bricken, B.~Chen, A.~Pearce, C.~Citro, E.~Ameisen, A.~Jones, et~al.
\newblock Scaling monosemanticity: Extracting interpretable features from Claude 3 Sonnet.
\newblock Anthropic, 2024.

\end{thebibliography}

\end{document}
